%%% FIELD proof-essential-feature-iff
Fix an enumerated saturated finite architecture \(T\), a deterministic table system in that architecture, and a point \(p_T\) in the product of source probability simplexes.  In particular, every source mass is nonnegative and the masses of each source sum to one.

First consider \(e=\operatorname{dir}(i,j)\).  By the definition of \(Z_e\), its finite sum can be indexed so that
\[
 Z_e(T,p_T)=\sum_{x\in\mathcal X_e} I_e(x)w_{p_T}(x),
 \qquad I_e(x)\in\{0,1\},
 \qquad w_{p_T}(x)\ge 0.
\]
Here \(x\) records an exogenous response context, \(w_{p_T}(x)\) is the corresponding monomial in the source masses and hence is the probability of that context, and \(I_e(x)=1\) exactly when there are two endogenous input assignments to the table of \(j\), differing only in the coordinate \(i\), at which the two outputs differ.  The comparison of endogenous inputs ranges over all table inputs, including assignments outside the observational support, by the definition of the graph-legality stratum; only the exogenous context is weighted by \(p_T\).

All summands in this finite sum are nonnegative.  Consequently,
\[
 Z_e(T,p_T)>0
 \quad\Longleftrightarrow\quad
 \text{there is an }x\in\mathcal X_e\text{ such that }
 I_e(x)=1\text{ and }w_{p_T}(x)>0.
\]
Indeed, a positive summand implies positivity of the sum, while positivity of a finite sum of nonnegative terms implies that at least one summand is positive.  Since \(w_{p_T}(x)>0\) means that \(x\) is a positive-probability exogenous context, the condition on the right is precisely actual functional dependence of the structural equation for \(j\) on the endogenous input \(i\), with dependence tested over all endogenous assignments.  By the definition of the minimal structural graph used in \(G_S\),
\[
 i\longrightarrow j\text{ is an edge of the minimal structural graph}
 \quad\Longleftrightarrow\quad
 Z_{\operatorname{dir}(i,j)}(T,p_T)>0.
\]

Now fix a nonempty \(C\subseteq W\), a vertex \(j\in C\), and \(e=\operatorname{src}(C,j)\).  Expanding its defining polynomial gives
\[
 Z_e(T,p_T)
 =\sum_{(u,u',x)\in\mathcal Y_e}
 J_e(u,u',x)\,p_C(u)p_C(u')w_{-C,p_T}(x),
\]
where \(u\ne u'\) are two slots of \(U_C\), \(x\) is a context of all other sources, \(w_{-C,p_T}(x)\ge0\) is its probability, and \(J_e(u,u',x)=1\) exactly when, at one common endogenous input assignment and the common other-source context \(x\), the table for \(j\) gives different outputs at \(U_C=u\) and \(U_C=u'\).  Again all terms are nonnegative, so
\[
 Z_e(T,p_T)>0
 \quad\Longleftrightarrow\quad
 \begin{gathered}
 \text{there exist }u\ne u'\text{ and }x\text{ with }
 p_C(u)>0,\ p_C(u')>0,\ w_{-C,p_T}(x)>0,\\
 \text{and }J_e(u,u',x)=1.
 \end{gathered}
\]
The right-hand side is exactly actual functional dependence of the equation for \(j\) on the exogenous source \(U_C\), outside exogenous null sets.  Hence
\[
 U_C\text{ is an actual parent of }j
 \quad\Longleftrightarrow\quad
 Z_{\operatorname{src}(C,j)}(T,p_T)>0.
\]
In particular, source values of zero mass cannot create an actual edge.

It follows that the complete positive-or-zero sign pattern indexed by \(\mathcal I_T\) recovers every directed edge and every source-to-child incidence of the minimal structural graph.  For each formal source \(U_C\), the definition
\[
 D_C(p_T)=\{j\in C:Z_{\operatorname{src}(C,j)}(T,p_T)>0\}
\]
therefore gives exactly its actual child set.  If \(D_C(p_T)=\varnothing\), that source has no actual child and may be discarded from the graph even when its mass vector is nondegenerate or its literal table entries vary only on null contexts.  No simplex-vertex condition is needed for this implication.  If distinct formal sources have the same nonempty realized child set \(D\), replacing them by their tuple source preserves their joint product law, lets each structural table unpack the same source coordinates, and has source-to-child incidences exactly \(D\).  Thus this merging preserves the induced SCM law, the normalized incidence pattern, and the latent projection, although it replaces several redundant latent vertices by one tuple vertex.  Discarding the empty realized sources and merging equal nonempty realized child sets therefore do not alter the subsequent selected-MAG projection.

The latent-and-selection projection is a deterministic finite operation on this normalized actual graph.  By the definition of \(\Psi_{T,G}\), it applies exactly that normalization and projection to the incidence pattern recovered above.  Therefore
\[
 G_S(M_{T,p_T})=G
 \quad\Longleftrightarrow\quad
 \Psi_{T,G}(p_T).
\]

Finally, a simultaneous permutation of source-value labels and their masses bijects the indices \(u,u'\) in the source-feature sums; a permutation of response labels bijects the corresponding response contexts; and a permutation of sources with the same realized child-set index permutes coordinates of the tuple source.  Each operation merely reindexes the displayed finite sums and leaves every equality \(Z_e=0\) and inequality \(Z_e>0\) unchanged after the corresponding label transport.  The normalized realized child sets and their projected selected MAG are consequently unchanged.  This proves all assertions.
